\documentclass[../numerical_methods.tex]{subfiles}
\begin{document}
\section{Aula 14 - 26 de Maio, 2026}
\subsection{Motivações}
\begin{itemize}
	\item Método de Gauss-Seidel
\end{itemize}
\subsection{Método de Gauss-Seidel}
O objetivo da aula de hoje é definir um processo iterativa tal que a sequência de vetores \(\{\underline{x}^{(0)}, \underline{x}^{(1)}, \underline{x}^{(2)}, \dotsc \}\) produzida por este processo \textit{convirja para a solução x}, independentemente da escolha do chute inicial \(\underline{x}^{(0)}\in \mathbb{R}^{n}\).

Mas pra que fazer isso? Muitas vezes, os métodos diretos podem ter problemas de implementação computacional; nestes casos, é necessário recorrer ao processo acima, que utiliza uma iteração convergente, mas aumentam o custo computacional.
\begin{def*}
	Uma sequência de vetores\(\{\underline{x}^{(0)}, \underline{x}^{(1)}, \underline{x}^{(2)}, \dotsc \}\) \textbf{convege para x} se
	\[
		\lim_{k\to \infty}\Vert \underline{x}^{(k)}-\underline{x} \Vert = 0. \; \square
	\]
\end{def*}
Denotamos o processo acima por \(\underline{x}^{(k)}\rightarrow \underline{x}.\)

A ideia principal do que faremos é criar um processo recursivo através de um sistema equivalente \(Ax = b\); dele, criaremos um sistema equivalente da forma
\[
	\underline{x} = C\underline{x} + \underline{g},
\]
onde \(C\in \mathbb{M}_{n}\) e \(\underline{g}\in \mathbb{R}^{n}\) são conhecidos. Daí, a partir de um chute inicial \(\underline{x}^{(0)}\), obtemos uma sequência \(\{\underline{x}^{(0)}, \underline{x}^{(1)}, \dotsc \}\) pelo processo iterativo
\[
	\boxed{\underline{x}^{(k+1)} = C \underline{x}^{(k)} + \underline{g}}, \quad k = 0, 1, 2, \dotsc
\]

Naturalmente, surgem algumas perguntas, principalmente:
\begin{itemize}
	\item É possível, dado \(A\underline{x}= \underline{b}\), obter um sistema equivalente da forma \(\underline{x} = C\underline{x} + \underline{g}\)? Sim! Basta tomar \(C = I-A \) e \(\underline{g} = \underline{b}\);
	\item Se a sequência \(\underline{x}^{(k)}\) converge para \(\underline{x}\), então \(\underline{x}\) será mesmo solução de \(A\underline{x} = \underline{b}\)? De fato, basta passar o limite em ambos os lados da equação recursiva, obtendo \(\underline{x} = C \underline{x} + \underline{g}\); pela hipótese de equivalência, segue que \(\underline{x} = \underline{x}\);
	\item Quando \(\underline{x}^{(k)}\rightarrow \underline{x}\) e quando parar o processo iterativo?
\end{itemize}
As duas primeiras puderam ser prontamente respondidas, mas as outras duas requerem mais cuidado.

\begin{def*}
	O \textbf{raio espectral} de uma matriz \(A\in \mathbb{M}_{n}\) é definido como
	\[
		\rho (A) = \max\limits_{i\in \{1, \dotsc , n\}} \{| \lambda_{i} |\},
	\]
	onde \(\lambda_{i}\) são os autovalores de A. \(\square\)
\end{def*}
Noutras palavras, o raio espectral é o maior autovalor de A, em módulo.

\hypertarget{general_criterion}{
	\begin{theorem*}[Critério Geral de Convergência]
		Seja \(\{\underline{x}^{(0)}, \underline{x}^{(1)}, \underline{x}^{(2)}, \dotsc \}\) a sequência gerada pelo processo iterativo
		\[
			\underline{x}^{(k+1)} = C \underline{x}^{(k)} + \underline{g}.
		\]
		Temos os critério de convergência:
		\begin{itemize}
			\item[1)] Se \(\Vert C \Vert_{M} < 1\), onde \(\Vert \cdot  \Vert_{M}\) é uma norma consistente qualquer, então a sequência converge;
			\item[2)] As iterações convergem para a solução se, e somente se, \(\rho (C) < 1\).
		\end{itemize}
	\end{theorem*}
}

Isso responde, em primeira vista, parte da última pergunta; quanto ao critério de parada, vamos tomar \(\varepsilon > 0\) e \(k_{\mathrm{max}}\in \mathbb{N}\) o natural correspondente ao passo máximo que faremos antes de para; então, temos alguns critérios:
\begin{itemize}
	\item \textbf{\underline{Erro absoluto}:} consiste em \(\Vert \underline{x}^{(k+1)} - \underline{x}^{(k)} \Vert <\varepsilon \);
	\item \textbf{\underline{Erro relativo}:} consiste no critério \(\displaystyle \frac{\Vert \underline{x}^{(k+1)} - \underline{x}^{(k)} \Vert}{\Vert \underline{x}^{(k+1)} \Vert} < \varepsilon \);
	\item \textbf{\underline{Teste de resíduo}:} corresponde ao teste \(\Vert \underline{b} - A\underline{x}^{(k)} \Vert < \varepsilon \);
	\item \textbf{\underline{Número máximo de iterações}:} simplesmente determinar \(k_{\mathrm{max}} = k\).
\end{itemize}

Dado \(Ax = b\) (e supondo que \(a_{ii}\neq 0\) sem perda de generalidade, para todo i), temos:
\[
	\left\{\begin{array}{ll}
		a_{11}x_1 + a_{12}x_2 + a_{13}x_3 + \dotsc + a_{1n}x_{n} = b_1 \\
		a_{21}x_1 + a_{22}x_2 + a_{23}x_3 + \dotsc + a_{2n}x_{n} = b_2 \\
		a_{31}x_1 + a_{32}x_2 + a_{33}x_3 + \dotsc + a_{3n}x_{n} = b_3 \\
		\vdots                                                         \\
		a_{n1}x_1 + a_{n2}x_2 + a_{n3}x_3 + \dotsc + a_{nn}x_{n} = b_n \\
	\end{array}\right.
\]
Este contexto será usado para introduzir o \textbf{método de Gauss-Jacobi}: ele transforma \(A\underline{x} = \underline{b}\) em \(\underline{x} = C\underline{x} + \underline{g}\) isolando cada coordenada i-ésima do vetor, ou seja, \(x_{i}\), na i-ésima equação do sistema. Logo,
Logo,
\[
	\left\{\begin{array}{ll}
		x_1 = (b_1-a_{12}x_2 - a_{13}x_3 -\dotsc - a_{1n}x_{n})/a_{11}       \\
		x_2 = (b_2-a_{21}x_1 - a_{23}x_3 -\dotsc - a_{2n}x_{n})/a_{22}       \\
		x_3 = (b_3-a_{31}x_1 - a_{32}x_2 -\dotsc - a_{3n}x_{n})/a_{33}       \\
		\vdots                                                               \\
		x_n = (b_n-a_{n2}x_2 - a_{n3}x_3 -\dotsc - a_{n, n-1}x_{n-1})/a_{nn} \\
	\end{array}\right.
\]
Logo, o sistema equivalente \(\underline{x} = C\underline{x} + \underline{g}\) é
\[
	C = \begin{bmatrix}
		0              & -a_{12}/a_{11} & -a_{13}/a_{11} & \cdots            & -a_{1n}/a_{11} \\
		-a_{21}/a_{22} & 0              & -a_{23}/a_{22} & \cdots            & -a_{2n}/a_{22} \\
		-a_{31}/a_{33} & -a_{32}/a_{33} & 0              & \cdots            & -a_{3n}/a_{33} \\
		\vdots         & \vdots         & \vdots         & \ddots            & \vdots         \\
		-a_{n1}/a_{nn} & -a_{n2}/a_{nn} & \cdots         & -a_{n,n-1}/a_{nn} & 0
	\end{bmatrix},\quad g=\begin{bmatrix}
		b_1/a_{11} \\
		b_2/a_{22} \\
		b_3/a_{33} \\
		\vdots     \\
		b_n/a_{nn}
	\end{bmatrix}.
\]
Portanto, dado o chute inicial \(\underline{x}^{(0)},\) o processo iterativo é dado por:
\[
	\left\{\begin{array}{ll}
		x_1^{(k+1)} = (b_1-a_{12}x_2^{(k)} - a_{13}x_3^{(k)} -\dotsc - a_{1n}     x_{n}^{(k)})/a_{11}  \\
		x_2^{(k+1)} = (b_2-a_{22}x_1^{(k)}- a_{23}x_3 ^{(k)}-\dotsc - a_{2n}      x_{n}^{(k)})/a_{22}  \\
		x_3^{(k+1)} = (b_3-a_{32}x_1^{(k)}- a_{33}x_2 ^{(k)}-\dotsc - a_{3n}      x_{n}^{(k)})/a_{33}  \\
		\vdots                                                                                         \\
		x_n^{(k+1)} = (b_n-a_{n2}x_2^{(k)} - a_{n3}x_3^{(k)} -\dotsc - a_{n, n-1}x_{n-1}^{(k)})/a_{nn} \\
	\end{array}\right.
\]
e o sistema fica \(\underline{x}^{(k+1)} = C\underline{x}^{(k)} + \underline{g}\):
\[
	C = \begin{bmatrix}
		0              & -a_{12}/a_{11} & -a_{13}/a_{11} & \cdots            & -a_{1n}/a_{11} \\
		-a_{21}/a_{22} & 0              & -a_{23}/a_{22} & \cdots            & -a_{2n}/a_{22} \\
		-a_{31}/a_{33} & -a_{32}/a_{33} & 0              & \cdots            & -a_{3n}/a_{33} \\
		\vdots         & \vdots         & \vdots         & \ddots            & \vdots         \\
		-a_{n1}/a_{nn} & -a_{n2}/a_{nn} & \cdots         & -a_{n,n-1}/a_{nn} & 0
	\end{bmatrix},\quad g=\begin{bmatrix}
		b_1/a_{11} \\
		b_2/a_{22} \\
		b_3/a_{33} \\
		\vdots     \\
		b_n/a_{nn}
	\end{bmatrix}.
\]
A forma como obter \(\underline{x}^{(k+1)} = C\underline{x}^{(k)} + \underline{g}\) a partir de \(Ax = b\) é considerando \(D\) como uma matriz diagonal formada pela diagonal de A, tal que
\[
	Ax = b \Longleftrightarrow (A-D+D)x = b \Longleftrightarrow (A-D)x + Dx = b.
\]
Dessa forma,
\[
	(A-D)\underline{x}^{(k)} + D\underline{x}^{(k+1)}=\underline{b} \Longleftrightarrow D\underline{x}^{(k+1)} = (D-A)\underline{x}^{(k)} + \underline{b}
\]
Portanto,
\[
	\underline{x}^{(k+1)} = \underbrace{(I-D^{-1}A)}_{C}\underline{x}^{(k)} + \underbrace{D^{-1}\underline{b}}_{\underline{g}}
\]
Implementando em python, ele fica:
\begin{listing}[H]
	\caption{Método de Gauss-Jacobi em Python}
	\begin{minted}[bgcolor = DarkBackground, linenos=true, breaklines]{Python}
import numpy as np

def gauss_jacobi(A, b, x0, tol, kmax=10000):
    A = A.astype(float)
    b = b.astype(float)
    d = np.diag(A)

    # Verifica se há zeros na diagonal para evitar divisão por zero
    if np.any(d == 0):
        raise ValueError("A diagonal principal contém zeros.")

    # Criando a matriz de iteração C e o vetor constante g
    # D_inv @ A é equivalente a dividir cada linha de A pelo elemento da diagonal
    D_inv_A = A / d[:, None]
    C = np.eye(len(b)) - D_inv_A
    g = b / d

    x = x0.astype(float)
    for k in range(kmax):
        x_novo = C @ x + g

        # Critério de parada pelo resíduo
        if np.linalg.norm(b - A @ x_novo) < tol:
            return x_novo, k + 1

        x = x_novo

    print('Erro: o método não converge no tempo estipulado.')
    return None, kmax
\end{minted}
\end{listing}
Podemos ir ao nosso primeiro critério de convergência: o método de Gauss-Jacobi converge para a solução de \(A\underline{x} = \underline{b}\), independentemente da escolha de \(\underline{x}^{(0)}\), se satisfizer um dos critérios:

\textbf{\underline{Critério das Linhas}:} o método converge se o máximo dos valores de \(\alpha_{k}\), onde
\[
	\alpha_{k} = \frac{\sum\limits_{\substack{j\neq k \\ j=1}}^{n}| a_{kj} |}{| a_{kk} |},
\]
for estritamente menor que 1.

\textbf{\underline{Critério das Colunas}:} o método converge se o máximo dos valores de \(\beta_{k}\), onde
\[
	\beta_{k} = \frac{\sum\limits_{\substack{i\neq k \\ i=1}}^{n}| a_{ik} |}{| a_{kk} |}.
\]
\begin{def*}
	Uma matriz que satisfaz o critério das linhas é chamada \textbf{estritamente diagonal dominante}. \(\square\)
\end{def*}
Quanto menor o valor de \(\alpha \), mais rápida será a convergência!

Para acelerar a convergência do método de Gauss-Jacobi, podemos usar os valores atualizados \(x_{1}^{(k+1)}, \dotsc , x_{i-1}^{(k+1)}\) e os restantes \(x_{i+1}^{(k)}, \dotsc , x_{n}^{(k)}\) no cálculo de \(x_{i}^{(k+1)}\):
\[
	\left\{\begin{array}{ll}
		x_1^{(k+1)} = (b_1-a_{12}x_2^{(k)} - a_{13}x_3^{(k)} -\dotsc - a_{1n}     x_{n}^{(k)})/a_{11}        \\
		x_2^{(k+1)} = (b_2-a_{22}x_1^{(k+1)}- a_{23}x_3 ^{(k)}-\dotsc - a_{2n}      x_{n}^{(k)})/a_{22}      \\
		x_3^{(k+1)} = (b_3-a_{32}x_1^{(k+1)}- a_{33}x_2 ^{(k+1)}-\dotsc - a_{3n}      x_{n}^{(k)})/a_{33}    \\
		\vdots                                                                                               \\
		x_n^{(k+1)} = (b_n-a_{n2}x_2^{(k+1)} - a_{n3}x_3^{(k+1)} -\dotsc - a_{n, n-1}x_{n-1}^{(k+1)})/a_{nn} \\
	\end{array}\right.
\]

Agora, outra forma de obtenção da forma matricial \(\underline{x}^{(k+1)} = C\underline{x}^{(k)} + g\), vamos utilizar o chamado \textbf{método de Gauss-Seidel}. Considere \(A = L + R\), sendo L a matriz triangular inferior de A e R a superior, mas sem a diagonal; desta forma,
\[
	A \underline{x}= \underline{b} \Longleftrightarrow (L+R)\underline{x} = \underline{b} \Longleftrightarrow L\underline{x}+R\underline{x} = \underline{b},
\]
tal que
\[
	L \underline{x}^{(k+1)}+R\underline{x}^{(k)} = \underline{b} \Longleftrightarrow L\underline{x}^{(k+1)} = - R\underline{x}^{(k)} + \underline{b}.
\]
Portanto, isolando o k+1-ésimo termo, obtemos
\[
	\boxed{\underline{x}^{(k+1)}=(-L^{-1}R)\underline{x}^{(k)}+L^{-1}\underline{b}}
\]
\begin{listing}[H]
	\caption{Método de Gauss-Seidel em Python}
	\begin{minted}[bgcolor = DarkBackground, linenos=true, breaklines]{Python}
import numpy as np

def gauss_seidel(A, b, x0, tol, kmax=10000):
    A = A.astype(float)
    b = b.astype(float)
    x = x0.astype(float)

    # Decomposição
    L = np.tril(A)      # Triangular inferior (com a diagonal)
    R = np.triu(A, 1)   # Triangular superior estrita

    for k in range(kmax):
        # Em vez de inv(L), resolvemos o sistema L @ x_novo = b - R @ x
        # Isso é equivalente a: x_novo = inv(L) @ (b - R @ x)
        rhs = b - R @ x
        x_novo = np.linalg.solve(L, rhs)

        # Critério de parada pelo resíduo
        if np.linalg.norm(b - A @ x_novo) < tol:
            return x_novo, k + 1

        x = x_novo

    print('Erro: o método não converge.')
    return None, kmax

# Exemplo
A = np.array([[4, -1, 0], [-1, 4, -1], [0, -1, 3]])
b = np.array([12, -1, 0])
x0 = np.zeros_like(b)

x, k = gauss_seidel(A, b, x0, 1e-6)
print(f"Solução: {x} | Iterações: {k}")
\end{minted}
\end{listing}

O método de Gauss-Seidel converge para a solução de \(A \underline{x} = \underline{b}\), independentemente da escolha do chute inicial, se ele satisfaz o critério de Sassenfeld: denotando por
\[
	\beta_1 = \frac{\sum\limits_{j=2}^{n}| a_{1j} |}{| a_{11} |} \quad\&\quad \beta_{i} = \frac{\sum\limits_{j=1}^{i-1}| a_{ij} |\beta_{j} + \sum\limits_{j=i+1}^{n}|a_{ij} |}{| a_{i} |},
\]
o método convergirá se
\[
	\beta =\max\limits_{1\leq i\leq n}\beta_{i} < 1,
\]
e quanto menor for o \(\beta \), mais rápida será a convergência.

A diferença entre os dois métodos é que o Gauss-Seidel converge mais rápido, mas o Gauss-Jacobi permite computações paralelas.

\subsection{Extra: Método dos Gradientes}
\begin{def*}
	Sejam \(A\in \mathbb{M}_{n},\; \underline{b}\in \mathbb{R}^{n}\) e c um número real. Uma \textbf{forma quadrática} é uma função \(F:\mathbb{R}^{n}\rightarrow \mathbb{R}\) escrita da forma
	\[
		F(x) = \frac{1}{2}\underline{x}^{T}A \underline{x} - \underline{x}^{T}\underline{b} + c. \; \square
	\]
\end{def*}

\begin{prop*}
	Se A é uma matriz simétrica, positiva e definida, então \(F(\underline{x})\) é minimizada pela solução de \(A \underline{x} = \underline{b}\).
\end{prop*}
\begin{proof*}[Esboço]
	O primeiro passo da prova é mostrar que \(\underline{x}\) é ponto crítico de F, ou seja,
	\[
		\nabla F(\underline{x}) = A \underline{x}-\underline{b} = 0 \Rightarrow A \underline{x} = \underline{b}.
	\]
	O segundo é mostrar que x é um mínimo, onde usa-se a \textit{matriz Hessiana}
	\[
		H = \biggl[\frac{\partial^{2}F(\underline{x})}{\partial x_{i}x_{j}^{}}\biggr] = A,
	\]
	e finalizar usando que A é SPD, portanto que x é ponto de mínimo. \qedsymbol
\end{proof*}

Com este contexto, o objetivo dos \textbf{métodos dos gradientes} é construir um processo iterativo \(\underline{x}^{(k)}\) (ou seja, uma sequência) que aproxime \(\underline{x}\). A ideia dele é descer o parabolóide da segunda forma no sentido contrário ao gradiente \(\nabla F(\underline{x})\), isto é,
\[
	- \nabla F(\underline{x}^{(k)}) = \underline{b} - A \underline{x}^{(k)} = \underline{r}^{(k)},
\]
onde \(\underline{r}^{(k)}\) é conhecido como \textbf{resíduo}. Dado um chute inicial \(\underline{x}^{(0)},\) obtemos \(\underline{x}^{(1)}\) da seguinte forma:
\[
	\underline{x}^{(1)}=\underline{x}^{(0)}-\alpha \nabla F(\underline{x}^{(0)}) = \underline{x}^{(0)} + \alpha \underline{r}^{(0)},
\]
tal que estaremos ``caminhando numa reta''... Mas, qual o tamanho do passo \(\alpha \)?

Para isso, o valor de \(\alpha \) tem que minimizar \(F(\underline{x})\) ao longo da reta \(\underline{x}^{(1)}=\underline{x}^{(0)} + \alpha \underline{r}^{(0)}\), que equivale a minimizar \(F(\underline{x}^{(1)})\). Logo,
\[
	\frac{\partial^{}F}{\partial \alpha ^{}}(\underline{x}^{(1)})=\nabla F(\underline{x}^{(1)}) \frac{\partial^{}\underline{x}^{(1)}}{\partial \alpha ^{}} = \overbrace{\nabla F(\underline{x}^{(1)})}^{=-\underline{r}^{(1)}}\cdot \underline{r}^{(0)}=0.
\]
Portanto, \(\underline{r}^{(1)}\) e \(\underline{r}^{(0)}\) são ortogonais (produto interno deles é nulo), e segue que
\begin{align*}
	0 = \underline{r}^{(1)}\cdot \underline{r}^{(0)} & =\biggl[\underline{b}-A\underline{x}^{(1)}\biggr]\cdot \underline{r}^{(0)}                                \\
	                                                 & =[\underline{b} - A(\underline{x}^{(0)} + \alpha \underline{r}^{(0)})]\cdot \underline{r}^{(0)}           \\
	                                                 & =[(\underline{b} - A \underline{x}^{(0)}) - \alpha A \underline{r}^{(0)}]\cdot \underline{r}^{(0)}        \\
	                                                 & =[\underline{r}^{(0)}-\alpha A \underline{r}^{(0)}]\cdot \underline{r}^{(0)}                              \\
	                                                 & =\underline{r}^{(0)}\cdot \underline{r}^{(0)} - \alpha (\underline{r}^{(0)} \cdot A \underline{r}^{(0)}).
\end{align*}
Daí,
\[
	\alpha = \frac{\underline{r}^{(0)}\cdot \underline{r}^{(0)}}{\underline{r}^{(0)}\cdot A \underline{r}^{(0)}},
\]
e o processo iterativo do método dos Gradientes é composto por:
\begin{align*}
	 & \underline{r}^{(k)} = \underline{b} - A \underline{x}^{(k)};                                                        \\
	 & \alpha ^{(k)} = \frac{\underline{r}^{(k)}\cdot \underline{r}^{(k)}}{\underline{r}^{(k)}\cdot A \underline{r}^{(k)}} \\
	 & \underline{x}^{(k+1)} = \underline{x}^{(k)} + \alpha^{(k)}\underline{r}^{(k)}.
\end{align*}
\begin{listing}[H]
	\caption{Método dos Gradientes em Python}
	\begin{minted}[bgcolor = DarkBackground, linenos=true, breaklines]{Python}
import numpy as np
from numpy import linalg as LA
def metodos_dos_gradientes(A, b, x, tol=1e-10, kmax=10000):
    # Garantir tipos float e converter para arrays caso sejam listas
    A = np.asanyarray(A).astype(float)
    b = np.asanyarray(b).astype(float)
    x = np.asanyarray(x).astype(float)

    # Verificação de simetria
    if not np.allclose(A, A.T):
        raise ValueError('A matriz A precisa ser simétrica.')

    # Verificação de Positiva Definida (via Cholesky - mais rápido que autovalores)
    try:
        np.linalg.cholesky(A)
    except np.linalg.linalg.LinAlgError:
        raise ValueError('A matriz A precisa ser positiva definida.')
    r = b - A @ x
    k = 0
    while LA.norm(r) > tol and k < kmax:
        p = r
        q = A @ p

        # Alpha é o tamanho do passo que minimiza f(x) na direção p
        alpha = (r @ r) / (p @ q)

        x = x + alpha * p
        r = b - A @ x  # Recalcular o resíduo para evitar acúmulo de erro numérico
        k += 1

    return x, k
\end{minted}
\end{listing}
\begin{listing}[H]
	\begin{minted}[bgcolor = DarkBackground, linenos=true, breaklines]{Python}

# Exemplo de uso
A = np.array([[3, 2], [2, 6]])
b = np.array([2, -8])
x0 = np.array([0, 0])

solucao, iteracoes = metodos_dos_gradientes(A, b, x0)
print(f"Solução: {solucao} em {iteracoes} iterações")
\end{minted}
\end{listing}

\end{document}
